
\begin{figure}[ht]
\begin{center}
\includegraphics[width=0.95\linewidth]{Fig/Result.png}

\caption{Qualitative results of all methods. (a) Input Image. (b) Ground Truth. (c) Grad-CAM \cite{selvaraju2017grad} (d) ScoreCAM \cite{wang2020score}. (e) LFI-CAM \cite{lee2021lfi}. (f) LayerCAM \cite{jiang2021layercam}. (g) Swin-MIL \cite{qian2022transformer}. (h) AME-CAM (ours). The image modalities of rows 1-4 are T1, T1-CE, T2, T2-FLAIR, respectively from the BraTS dataset.
}
\label{Fig.Result}
\end{center}
\end{figure}

\subsection{Quantitative and Qualitative Comparison with State-of-the-art}
\label{result:cam_sota}

In this section, we compare the segmentation performance of the proposed AME-CAM with five state-of-the-art weakly-supervised segmentation methods, namely Grad-CAM \cite{selvaraju2017grad}, ScoreCAM \cite{wang2020score}, LFI-CAM \cite{lee2021lfi}, LayerCAM \cite{jiang2021layercam}, and Swin-MIL \cite{qian2022transformer}. We also compare with an unsupervised approach C\&F \cite{chen2020medical}, the supervised version of C\&F, and the supervised Optimized U-net \cite{futrega2021optimized} to show the comparison with non-CAM-based methods. We acknowledge that the results from fully supervised and unsupervised methods are not directly comparable to the weakly supervised CAM methods. Nonetheless, these methods serve as interesting references for the potential performance ceiling and floor of all the CAM methods.

Quantitatively, Grad-CAM and ScoreCAM result in low dice scores, demonstrating that they have difficulty extracting the activation of medical images. LFI-CAM and LayerCAM improve the dice score in all modalities, except LFI-CAM in T1-CE and T2-FLAIR. Finally, the proposed AME-CAM achieves optimal performance in all modalities of the BraTS dataset.

Compared to the unsupervised baseline (UL), C\&F is unable to separate the tumor and the surrounding tissue due to low contrast, resulting in low dice scores in all experiments. With pixel-wise labels, the dice of supervised C\&F improves significantly. Without any pixel-wise label, the proposed AME-CAM outperforms supervised C\&F in all modalities.

The fully supervised (FSL) Optimized U-net achieves the highest dice score and IoU score in all experiments. However, even under different levels of supervision, there is still a performance gap between the weakly supervised CAM methods and the fully supervised state-of-the-art. This indicates that there is still potential room for WSSS methods to improve in the future.

Qualitatively, Fig. \ref{Fig.Result} shows the visualization of the CAM and segmentation results from all six CAM-based approaches under four different modalities from the BraTS dataset. Grad-CAM (Fig. \ref{Fig.Result}(c)) results in large false activation region, where the segmentation mask is totally meaningless. ScoreCAM eliminates false activation corresponding to air. LFI-CAM focus on the exact tumor area only in the T1 and T2 MRI (row 1 and 3). Swin-MIL can hardly capture the tumor region of the MRI image, where the activation is noisy. Among all, only LayerCAM and the proposed AME-CAM successfully focus on the exact tumor area, but AME-CAM reduces the under-estimation of the tumor area. This is attributed to the benefit provided by aggregating activation maps from different resolutions.


\subsection{Ablation Study}



\begin{table}
\centering
\caption{Ablation study for aggregation phase using T1 MRI images from the BraTS dataset. Avg. ME denotes that we directly average four activation maps generated by the multiple-exit phase. The dice score, IoU, and the HD95 are reported in the form of mean$\pm$std.}
\label{Table:ablation_aggregation}
\small
\begin{tabular}{c|c|c|l}
\hline
Method                            & Dice $\uparrow$        & IoU $\uparrow$         & \multicolumn{1}{c}{HD95 $\downarrow$}           \\ \hline
Avg. ME                           & 0.617$\pm$0.121 & 0.457$\pm$0.121 & 23.603$\pm$20.572  \\
Avg. ME+C$^2$AM\cite{xie2022c2am} & 0.484$\pm$0.256 & 0.354$\pm$0.207 & 69.242$\pm$121.163 \\
AME-CAM (ours)                    & \textbf{0.631$\pm$0.119} & \textbf{0.471$\pm$0.119} & \textbf{21.813$\pm$18.219} \\ \hline
\end{tabular}
\end{table}
\textbf{Effect of Different Aggregation Approaches:} In Table \ref{Table:ablation_aggregation}, we conducted an ablation study to investigate the impact of using different aggregation approaches after extracting activations from the multiple-exit network. We aim to demonstrate the superiority of the proposed attention-based aggregation approach for segmenting tumor regions in T1 MRI of the BraTS dataset. Note that we only report the results for T1 MRI in the BraTS dataset. Please refer to the supplementary material for the full set of experiments.

As a baseline, we first conducted the average of four activation maps generated by the multiple-level activation extraction (Avg. ME). We then applied C$^2$AM \cite{xie2022c2am}, a state-of-the-art CAM-based refinement approach, to refine the result of the baseline, which we call "Avg. ME+C$^2$AM". However, we observed that C$^2$AM tended to segment the brain region instead of the tumor region due to the larger contrast between the brain tissue and the air than that between the tumor region and its surrounding tissue. Any incorrect activation of C$^2$AM also led to inferior results, resulting in a degradation of the average dice score from 0.617 to 0.484. In contrast, the proposed attention-based approach provided a significant weighting solution that led to optimal performance in all cases.

\begin{table}[h]
\centering
% \setlength{\tabcolsep}{3pt}
\caption{Ablation study for using single-exit from $M_1$, $M_2$, $M_3$ or $M_4$ of Fig. \ref{Fig.AME-CAM} and the multiple-exit using results from $M_2$ and $M_3$ and using all exits (AME-CAM). The experiments are done on the T1-CE MRI of BraTS dataset. The dice score, IoU, and the HD95 are reported in the form of mean$\pm$std.}
\label{Table:ablation_single_exit}
\begin{tabular}{c|c|c|c|l}
\hline
\multicolumn{2}{c|}{Selected Exit}                & Dice $\uparrow$        & IoU $\uparrow$         & \multicolumn{1}{c}{HD95 $\downarrow$}           \\ \hline
\multirow{4}{*}{Single-exit}   & $M_1$          & 0.144$\pm$0.184 & 0.090$\pm$0.130 & 74.249$\pm$62.669           \\
              & $M_2$          & 0.500$\pm$0.231 & 0.363$\pm$0.196 & 43.762$\pm$85.703           \\
              & $M_3$          & 0.520$\pm$0.163 & 0.367$\pm$0.141 & 43.749$\pm$54.907           \\
              & $M_4$          & 0.154$\pm$0.101 & 0.087$\pm$0.065 & 120.779$\pm$44.548          \\ \hline
\multirow{2}{*}{Multiple-exit} & $M_2+M_3$      & 0.566$\pm$0.207 & 0.421$\pm$0.186 & 27.972$\pm$56.591           \\
              & AME-CAM (ours) & \textbf{0.695$\pm$0.095} & \textbf{0.540$\pm$0.108} & \textbf{18.129$\pm$12.335}           \\ \hline
\end{tabular}
\end{table}
\textbf{Effect of Single-Exit and Multiple-Exit:} Table \ref{Table:ablation_single_exit} summarizes the performance of using single-exit from $M_1$, $M_2$, $M_3$, or $M_4$ of Fig. \ref{Fig.AME-CAM} and the multiple-exit using results from $M_2$ and $M_3$, and using all exits (AME-CAM) on T1-CE MRI in the BraTS dataset.

The comparisons show that the activation map obtained from the shallow layer $M_1$ and the deepest layer $M_4$ result in low dice scores, around 0.15. This is because the network is not deep enough to learn the tumor region in the shallow layer, and the resolution of the activation map obtained from the deepest layer is too low to contain sufficient information to make a clear boundary for the tumor. Results of the internal classifiers from the middle of the network ($M_2$ and $M_3$) achieve the highest dice score and IoU, both of which are around 0.5.

To evaluate whether using results from all internal classifiers leads to the highest performance, we further apply the proposed method to the two internal classifiers with the highest dice scores, i.e., $M_2$ and $M_3$, called $M_2+M_3$. Compared with using all internal classifiers ($M_1$ to $M_4$), $M_2+M_3$ results in 18.6\% and 22.1\% lower dice and IoU, respectively. In conclusion, our AME-CAM still achieves the optimal performance among all the experiments of single-exit and multiple-exit.

Other ablation studies are presented in the supplementary material due to space limitations.